\section{Conclusions}\label{sec:conclusion}

This dissertation asked what a fiscal authority can learn from LLM-based exposure measures, given that they are demonstrably model-dependent. The answer is that it can learn the difference but not the level. Where the question is posed as a gap across regions, and the gap is read scale-free, the measure holds; levels carry a systematic bias that averaging cannot remove and will not support the analyses built on top of them.

Scotland sits about two percentile points below the rest of the UK in the UK occupational exposure distribution, and 0.86 index points on the baseline model's 0--100 scale. The composed 95\% interval on that figure, $[-1.07, -0.69]$, has most of the width coming from survey sampling, and describes precision within one model rather than agreement across models. London accounts for about two-fifths of the gap, and Scotland ranks fifth of the twelve UK nations and regions; against the ex-London comparator the shortfall narrows to 0.51 index points, so the finding is less that Scotland is insulated than that London is unusually exposed. Behind the gap is a composition underweight in professional-management work, which scores high, and overweight in the care sector, which scores low, both of which persist against the ex-London benchmark.

The robustness of the differential is established mathematically and then verified through a re-scoring exercise. Proposition 1 shows the differential removes the common component of any model-specific bias, with Corollary 1 bounding what survives. Corollary 2 covers the case the data in fact present, where the alternative models compress the scale rather than reordering occupations, and since Proposition 1 does not protect against this, a standardised or rank-based differential is supported. Specifically, I find on the re-score that the raw gap roughly halves while absolute levels move by up to 9.1 index points, three-quarters to nine-tenths of that movement being the affine rescaling of Corollary 2. The remainder is small in the terms Corollary 1 sets: the model bias correlates with the employment-share gap at $-0.23$ and $-0.24$, so the non-affine residual reaches only about a quarter of the Cauchy--Schwarz worst case. The sign of the differential, its composition, and its scale-free magnitude are all robust to the choice of scoring model.

Carried into productivity projections à la \citet{acemoglu2025}, the gap implies a ten-year TFP differential of $-0.031$ percentage points, or three basis points, which translates into a net-position sensitivity of about \pounds\fiscCen~million a year. Against the \pounds969 million forecast net position, this figure is modest and well inside the SFC's own income tax forecast-error ranges, but it is persistent, and applies afresh at each three-year reconciliation.

Since it is a theme of this paper not to trust levels, it would be remiss to report a fiscal translation figure and abstract from model dependency. Every step from score index to the net position multiplies by a coefficient common to the two regions, thus by Corollary 3 making the chain a linear functional: the Proposition 1 cancellation reaches the net position, and its uncertainty can be priced by rescaling. Doing so orders the sources. Scoring noise moves the annual figure least, by \fiscNoisePct~per cent either side of the centre; survey sampling and the adoption scenario shift it by around two-fifths of the central figure each; and the choice of scoring LLM by \fiscModelPct~per cent, which averaging cannot reduce. The wage gradient makes the same point at worker level, where the errors-in-variables correction is 1.6\%, changing nothing.

For a forecasting body the guidance is short. Use continuous indices in place of classified shares, because a classifier that codes 91\% of employment as exposed censors any variation required for regional comparison. Trust differentials over levels, as differencing removes the bias that makes levels a property of the scoring model rather than of the labour market. And expect the choice of model to move magnitudes. Design around it: report the sign and the percentile statement, neither of which inherits the model's scaling, and let direction and persistence take precedence over the index-point figure.

The limitations are as follows. Since no instrument for AI exposure yet matches the commuting-zone design of the automation literature, nothing here is causally identified, and so what I report are compositional measurements rather than displacement effects. Assumption 1, region invariance of within-occupation task content, cannot be tested without UK regional task data that no source provides, though the industry-margin diagnostics of \cref{app:data} bound the strain it is under. 
%The productivity projection imports the Acemoglu calibration wholesale, so its levels inherit whatever that calibration gets wrong, which is why I rest on the differential in the implied paths. 
Lastly, the fiscal translation is merely an accounting exercise using revenue elasticity computed from the Scottish schedule rather than assumed, with the pass-through from productivity to earnings imposed, and the population channel running through per-capita indexation is left out entirely.

What the framework offers extends beyond this application. Any exposure-based
regional estimand inherits the same common-mode structure, for any pair of regions
and any scoring model: an employment-weighted contrast whose weights sum to one
differences out whatever the model does uniformly, and what remains is bounded by
the alignment of the model's bias with the share gap. 
%The calibration is cheap to re-run as the models change, which they will. 
And where magnitude is not immaterial as it is here, the choice of scoring model can impact the figures institutions lean on even more than the data do.

The natural extensions centre around data this
project could not obtain. Employer-reported ASHE earnings under a secure Data Access Agreement would replace recalled pay and supply an external check on the pooled occupational shares. The secure-access APS would support the four-digit
differential the public file compresses, along with unit-group employment weights
for the second step of the crosswalk. 
%And a longer panel would support the first-difference specification the current window cannot identify.